\documentclass[10pt]{article}
\usepackage[T1]{fontenc}
\usepackage[utf8]{inputenc}
\usepackage[french]{babel}
\usepackage{amsmath,amssymb,mathrsfs,mathpazo}
\usepackage{geometry}
\geometry{a4paper,margin=2cm}
\usepackage{graphicx,array,tabularx,tabularray,colortbl,multirow,multicol,booktabs,enumitem,pifont,ulem,fancybox,diagbox,cancel}
\usepackage{pstricks,pstricks-add,pst-plot,pst-tree,pst-node,pst-text,pst-func,pst-math,pst-bspline,pst-3dplot,pst-eucl,pst-solides3d}
\let\pscalculate\undefined
\usepackage{pst-bezier}
\usepackage{tikz,pgfplots,tkz-tab}
\usetikzlibrary{arrows,arrows.meta,patterns,calc,positioning,trees,decorations.markings,shapes.geometric,shapes.misc,shapes.arrows,intersections,decorations.pathmorphing,decorations.pathreplacing,backgrounds,shadows,babel}
\pgfplotsset{compat=1.18}
\usepgfplotslibrary{fillbetween}
\usepackage[np]{numprint}
\usepackage{listings,makecell,pas-tableur,twoopt}
\let\barre\undefined
\usepackage[pstricks]{tablvar}
\definecolor{darkgreen}{RGB}{0,100,0}
\providecommand{\red}{\color{red}}
\providecommand{\blue}{\color{blue}}
\providecommand{\green}{\color{green}}
\providecommand{\black}{\color{black}}
\providecommand{\white}{\color{white}}
\providecommand{\gray}{\color{gray}}
\providecommand{\cyan}{\color{cyan}}
\providecommand{\magenta}{\color{magenta}}
\providecommand{\yellow}{\color{yellow}}
\providecommand{\violet}{\color{violet}}
\providecommand{\orange}{\color{orange}}
\providecommand{\euro}{\texteuro}
\providecommand{\EUR}{\texteuro}
\providecommand{\ieme}{\textsuperscript{e}}
\providecommand{\degres}{\ensuremath{^\circ}}
\providecommand{\degree}{\ensuremath{^\circ}}
\providecommand{\vv}[1]{\overrightarrow{#1}}
\providecommand{\Coord}[1]{#1}
\providecommand{\decosix}{}
\providecommand{\starredbullet}{\ensuremath{\bullet}}
\pagestyle{empty}
\setlength{\parindent}{0pt}
\usepackage{adjustbox}\DeclareUnicodeCharacter{2236}{:}\DeclareUnicodeCharacter{2190}{\ensuremath{\leftarrow}}
\providecommand{\R}{\mathbb{R}}
\providecommand{\N}{\mathbb{N}}
\providecommand{\Z}{\mathbb{Z}}
\providecommand{\Q}{\mathbb{Q}}
\providecommand{\D}{\mathbb{D}}
\providecommand{\C}{\mathbb{C}}
\providecommand{\e}{{\mathrm{e}}}
\providecommand{\vect}[1]{\overrightarrow{#1}}
\providecommand{\vectt}[1]{\overrightarrow{\mathrm{#1}}}
\providecommand{\vvt}[1]{\overrightarrow{\mathrm{#1}}}
\providecommand{\np}[2][]{{#2~#1}}
\providecommand{\vv}[1]{{\overrightarrow{#1}}}
\renewcommand{\d}{{\mathrm{d}}}
\providecommand{\degres}{{^\circ}}
\providecommand{\degre}{{^\circ}}
\providecommand{\blue}{\color{blue}}
\providecommand{\red}{\color{red}}
\providecommand{\enskip}{\quad}
\providecommand{\llbracket}{\lbrack\!\lbrack}
\providecommand{\rrbracket}{\rbrack\!\rbrack}
\providecommand{\euro}{{\text{€}}}
\providecommand{\no}{{\text{n°}}}
\providecommand{\interval}[2]{{[#1;#2]}}
\providecommand{\intervFO}[2]{{[#1;#2[}}
\providecommand{\dotfill}{\ldots}

\providecommand{\ds}{\displaystyle}
\providecommand{\textding}[1]{#1}
\providecommand{\Oij}{\(\mathrm{O};\vec{\imath},\vec{\jmath}\)}
\providecommand{\Oijk}{\(\mathrm{O};\vec{\imath},\vec{\jmath},\vec{k}\)}
\providecommand{\Ouv}{\(\mathrm{O};\vec{u},\vec{v}\)}
\providecommand{\intervFF}[2]{\ensuremath{[#1;#2]}}
\renewcommand{\euro}{\ifmmode\text{\texteuro}\else\texteuro\fi}
\DeclareUnicodeCharacter{00B0}{\ifmmode^{\circ}\else\textdegree\fi}
\DeclareUnicodeCharacter{20AC}{\ifmmode\text{\texteuro}\else\texteuro\fi}
% Énoncés : Éducation nationale. Transcription et corrigés : APMEP, Denis Vergès et contributeurs. Découpage et classement : CoopMaths, Sébastien Lozano, Éric Elter et contributeurs. Adaptation de la présentation : Maths Explorer.
% CC BY-SA : https://coopmaths.fr/www/about/
% https://coopmaths.fr/alea/static/bac/2023/tex/bac_2023_09_sujet1_polynesie_3_cor.tex
\begin{document}
\noindent\textit{APMEP / CoopMaths. CC BY-SA. Présentation adaptée par Maths Explorer.}\par\medskip
\medskip

On considère la fonction $f$ définie sur $\R$ par 
$f(x) = \dfrac34 x^2 - 2x + 3.$

\begin{enumerate}
\item 
$f'(x)=\dfrac{3}{2}x -2 = \dfrac{3}{2}\left (x-\dfrac{4}{3}\right )$, donc $f'(x)$ est du signe de $x-\dfrac{4}{3}$ donc s'annule et change de signe pour $x=\dfrac{4}{3}$.

$f$ est une fonction polynôme donc:

$\ds\lim_{x\to -\infty} f(x) = \ds\lim_{x\to -\infty} \dfrac{3}{4}x^2 = +\infty$
et
$\ds\lim_{x\to +\infty} f(x) = \ds\lim_{x\to +\infty} \dfrac{3}{4}x^2 = +\infty$.

De plus, $f\left(\dfrac43\right) = \dfrac34 \times \dfrac43 \times \dfrac43 - 2\times \dfrac43 + 3 = \dfrac43 - \dfrac83 + \dfrac93 = \dfrac53$.

On établit le tableau des variations de $f$:

\begin{center}
{\renewcommand{\arraystretch}{1.3}
\psset{nodesep=3pt,arrowsize=2pt 3}  
\def\esp{\hspace*{1.5cm}}
\def\hauteur{0pt}
$\begin{array}{|c| *4{c} c|}
\hline
 x & -\infty & \esp & \frac{4}{3} & \esp & +\infty \\
\hline
x-\frac{4}{3} &  &  \pmb{-} & \vline\hspace{-2.7pt}0 & \pmb{+} & \\
 \hline
f'(x) &  &  \pmb{-} & \vline\hspace{-2.7pt}0 & \pmb{+} & \\  
\hline
  & \Rnode{max1}{+\infty}  &  &  &  & \Rnode{max2}{+\infty}   \\
f(x) & &  & & &  \rule{0pt}{\hauteur} \\
 &  & &   \Rnode{min}{\frac{5}{3}} & & \rule{0pt}{\hauteur}
\ncline{->}{max1}{min} \ncline{->}{min}{max2}\\
\hline
\end{array}$
}
\end{center}

\item 
D'après le résultat précédent la fonction $f$ est croissante sur l'intervalle $\left[\dfrac43~;~2\right]$.

On a donc : $\dfrac43 \leqslant x \leqslant 2 \Longrightarrow f\left(\dfrac43\right) \leqslant f(x) \leqslant f(2)$.

Or $f\left(\dfrac43\right) = \dfrac53$ et $f(2) = \dfrac34 \times 4 - 2 \times 2 + 3 = 3 - 4 + 3 = 2$.

Donc $\dfrac53 \leqslant f(x) \leqslant  2$ et \emph{a fortiori} : $\dfrac43 \leqslant f(x) \leqslant 2$.

\item




$f(x)-x = \dfrac{3}{4} x^2 -2x+3-x = \dfrac{3}{4}x^2 - 3x +3 = \dfrac{3}{4}\left ( x^2 -4x+4 \right ) = \dfrac{3}{4} \left (x-2\right )^2$ 

Pour tout réel $x$, $(x-2)^2\geqslant 0$ donc $f(x)-x \geqslant 0$ donc $x\leqslant f(x)$.
\end{enumerate}

On considère la suite $\left(u_n\right)$ définie par un réel $u_0$ et pour tout  $n\in\N$ : $u_{n+1} = f\left(u_n\right).$

On a donc, pour tout entier naturel $n$,
$u_{n+1} = \dfrac34 u_n^2 - 2u_n + 3.$

\begin{enumerate}[resume]
\item Étude du cas: $\dfrac43 \leqslant u_0 \leqslant  2$.
	\begin{enumerate}
		\item On va démontrer par récurrence que  la propriété $u_n \leqslant u_{n+1} \leqslant 2$ est vraie pour tout entier naturel $n$.

\begin{list}{\textbullet}{}
\item \textbf{Initialisation}

D'après la question précédente, pour $x\in\left [ \dfrac{4}{3}\;;\;2\right ]$, on a $x\leqslant f(x)$. \\
Or $u_0\in\left [ \dfrac{4}{3}\;;\;2\right ]$ donc $u_0\leqslant f(u_0)$, ce qui revient à $u_0\leqslant u_1$.

De plus, si $x\in\left [ \dfrac{4}{3}\;;\;2\right ]$, $f(x)\in\left [ \dfrac{4}{3}\;;\;2\right ]$.\\
Or $u_0\in\left [ \dfrac{4}{3}\;;\;2\right ]$ donc  $f(u_0)\in\left [ \dfrac{4}{3}\;;\;2\right ]$ et donc $u_1 \in\left [ \dfrac{4}{3}\;;\;2\right ]$ et donc $u_{n+1}\leqslant 2$.

Donc $u_0\leqslant u_1 \leqslant 2$; la propriété est vraie au rang $n=0$.

\item \textbf{Hérédité}

On suppose la propriété vraie au rang $n$, c'est-à-dire $u_n\leqslant u_{n+1}\leqslant 2$.

On est dans l'intervalle $\left [ \dfrac{4}{3}\;;\;2\right ]$ donc la fonction $f$ est croissante; on en déduit: $f(u_n)\leqslant f(u_{n+1})\leqslant f(2)$.

Or $f(u_n)=u_{n+1}$, $f(u_{n+1})=u_{n+2}$ et $f(2)=2$.

Donc on a: $u_{n+1}\leqslant u_{n+2}\leqslant 2$; la propriété est héréditaire.
\item \textbf{Conclusion}

La propriété est vraie au rang 0 et elle est héréditaire pour tout $n\geqslant 0$; d'après le principe de récurrence, elle est vraie pour tout  $n\in\N$.
\end{list}
		\item 
Pour tout $n$, on a $u_n \leqslant u_{n+1}\leqslant 2$ donc la suite $\left(u_n\right)$ est croissante et majorée; d'après le théorème de la convergence monotone, la suite 
$\left(u_n\right)$ est convergente vers une limite $\ell$ telle que $\ell \leqslant 2$.

		\item 
La suite $\left(u_n\right)$ est définie par $f(u_n)=u_{n+1}$ où $f$ est une fonction polynôme donc continue. 

La suite $\left(u_n\right)$ est convergente vers $\ell$ donc la limite $\ell$ vérifie $f(\ell) = \ell$.

$f(\ell) = \ell \iff f(\ell)-\ell=0 \iff \dfrac{3}{4}\left(\ell -2\right )^2= 0 \iff \ell = 2$.

La suite $\left(u_n\right)$ converge donc vers 2.


	\end{enumerate}
\item Étude du cas: $u_0 = 3$.
On admet que dans ce cas la suite $\left(u_n\right)$ tend vers $+\infty$.

On complète la fonction \og \texttt{seuil} \fg{} suivante écrite en Python, afin qu'elle renvoie la plus petite valeur de $n$ telle que $u_n$ soit supérieur ou égal à 100.

\begin{center}
\texttt{
\begin{tabular}{|l|}\hline
def seuil() :\\
\quad  u = 3\\
\quad  n = 0\\
\quad while {\blue u < 100}\\
\quad\qquad  u = {\blue 3*u*u/4 - 2*u + 3}\\
\quad\qquad  n = {\blue n + 1}\\
\quad return n\\ \hline
\end{tabular}}
\end{center}


\item Étude du cas : $u_0 > 2$.


Pour tout réel $x$, on a: $x\leqslant f(x)$ donc, pour tout $n$, $u_n\leqslant f(u_n)$ soit $u_n\leqslant u_{n+1}$ ; la suite $(u_n)$ est donc croissante.

On en déduit que pour tout $n$, on a: $u_n\geqslant u_0$.

Si la suite $(u_n)$ converge vers un réel $\ell$, on aura donc: $\ell \geqslant u_0$.

On a vu que la seule limite possible de la suite $(u_n)$ était $\ell = 2$; donc on ne peut pas avoir $\ell \geqslant u_0$ car $u_0>2$.

Pour $u_0>2$, la suite $(u_n)$ n'est donc pas convergente.

\end{enumerate}

\bigskip

\end{document}
